The saints whom the Church celebrates together have reached the happiness
Christ promises on the mountain. The people celebrating them still ask for
mercy, release from sin, and protection. All Saints brings those two conditions
into the same act of worship. The Gospel names the purity and peace of God's
children, including their fidelity under persecution. The prayers ask God to
give the worshippers the saving gifts they need, with the saints praying for
them. The distinction between the saints' happiness and the worshippers' need
explains why a feast of so much joy contains so much petition.

John's vision also requires patient attention. Before the multitude from every
nation appears, John hears an exact number and a list of twelve tribes. Bede
explains how such a definite count can signify the Church in its completeness.
Augustine explains why death does not remove Christ's faithful members from
his Church. Together their arguments help us understand the celebration of one
people whose members praise God in different conditions. A second
interpretation concerns the life Christ teaches those people to live.
Chrysostom connects the refreshment promised to the burdened with Christ's
humility; Augustine explains the Beatitudes as the formation of desire,
purity, and peace. Their accounts address the apparent difficulty of finding
rest while hungering for justice and enduring persecution.

On Sunday, 1 November 2026, the first-class feast of All Saints is celebrated
with a commemoration of the Twenty-third Sunday after Pentecost. The Sunday
adds a second Collect, Secret, and Postcommunion. The feast's Lesson, Gospel,
and chants remain those of All Saints; the Common Preface is said. The Creed
is expressly appointed.\footnote{\work{Missale Romanum}, Vatican editio
typica 1962, pp.~702--703; \work{Rubricae generales} 91, 106, 108--109,
111, 113, 120(d); \work{Rubricae generales Missalis romani} 475, 480,
482--483, 494(b), 498, 505. The three Sunday prayers are on pp.~416--417.}
Those additional prayers make the congregation's present need unusually
explicit: sins have entangled God's people through their frailty; a gift
already received needs completion; participation in divine things calls for
protection amid human dangers.

\section{The propers at a glance}
\begin{maptable}
Introit & \work{Gaudeamus}; Ps.~32:1 (33:1) & The Church rejoices in the Lord with angels praising the Son.\\
Collect & \work{Omnipotens sempiterne Deus} & One celebration honors all the saints; many intercessors accompany the request for mercy.\\
Sunday Collect & \work{Absolve, quaesumus} & God's people ask release from sins contracted through frailty.\\
Lesson & Apoc.~7:2--12 & The sealed tribes and a multitude from every nation worship God and the Lamb.\\
Gradual & Ps.~33:10,11b (34:10,11b) & The saints fear God and seek him as the giver of good.\\
Alleluia & Matt.~11:28 & Christ invites the burdened to receive refreshment.\\
Gospel & Matt.~5:1--12a & Jesus promises the kingdom to the humble, merciful, pure, peaceful, and faithfully persecuted.\\
Offertory & Wis.~3:1--3, selected & The just who appear to die remain in God's custody and peace.\\
Secret & \work{Munera tibi, Domine} & The offerings honor all the just and ask saving benefit for the offerers.\\
Sunday Secret & \work{Pro nostrae servitutis} & The sacrifice of praise asks growth of service and completion of an unmerited gift.\\
Common Preface & \work{Vere dignum} & Through Christ, human thanksgiving asks admission among the heavenly praises.\\
Communion & Matt.~5:8--10 & Purity, peacemaking, and persecution for justice are repeated to the communicants.\\
Postcommunion & \work{Da, quaesumus} & The faithful ask lasting joy in the saints and their continual supplication.\\
Sunday Postcommunion & \work{Quaesumus, omnipotens} & Those admitted to divine participation ask protection from human dangers.\\
\end{maptable}

Psalm references give Vulgate numbering first and the common modern number in
parentheses. The Lesson's traditional name, Apocalypse, and Revelation name
the same biblical book. The four senses developed below distinguish what the
texts say in context, their fulfillment in Christ and his body, the conversion
they teach, and their final fulfillment in the life to come.

\tableofcontents
